Tại phần này, bạn sẽ thấy bản tóm tắt đề xuất dự án phát triển hệ thống
Y tế số, bao gồm mục tiêu, kiến trúc hạ tầng AWS Cloud, luồng nghiệp vụ
và ước tính ngân sách vận hành.

\vspace{0.5em}\noindent\textbf{Smart Healthcare AI Triage \& Appointment Booking Platform}

\vspace{0.5em}\noindent\textbf{Giải pháp Cloud-native hỗ trợ Sàng lọc Bệnh Y khoa qua AI và
Quản lý Lịch hẹn Trực tuyến}

\vspace{0.5em}\noindent\textbf{1. Tóm tắt điều hành}

Smart Healthcare Platform được thiết kế nhằm hiện đại hóa quy trình tiếp
nhận, sàng lọc ban đầu và đặt lịch khám bệnh tại các phòng khám/bệnh
viện. Hệ thống ứng dụng mô hình trí tuệ nhân tạo (AI) trên
\textbf{Amazon SageMaker} để phân loại triệu chứng của bệnh nhân theo 3
cấp độ (Đỏ - Vàng - Xanh) và tự động tạo báo cáo tóm tắt chuẩn y khoa
cho bác sĩ trước ca khám. Nền tảng được xây dựng theo kiến trúc
Microservices/Cloud-native trên nền tảng \textbf{AWS Cloud} (Cognito,
EC2, RDS PostgreSQL, S3, CloudWatch), phục vụ hàng nghìn lượt truy cập
đồng thời từ bệnh nhân, bác sĩ và đội ngũ lễ tân.

\vspace{0.5em}\noindent\textbf{2. Tuyên bố vấn đề}

\emph{Vấn đề hiện tại}\\
Các cơ sở y tế hiện nay thường gặp tình trạng quá tải tại quầy tiếp đón.
Quy trình đặt lịch hẹn truyền thống qua điện thoại hoặc tại quầy dễ gây
ra tình trạng trùng lịch (Double-booking), thời gian chờ đợi của bệnh
nhân kéo dài. Bên cạnh đó, bác sĩ tốn nhiều thời gian hỏi lại từ đầu các
triệu chứng cơ bản của bệnh nhân do thiếu thông tin chuẩn bị trước
(Pre-consultation summary).

\emph{Giải pháp}\\
Nền tảng cung cấp giải pháp toàn diện bao gồm:\\
1. \textbf{Chatbot AI thời gian thực:} Giao tiếp với bệnh nhân qua
WebSocket, thu thập triệu chứng và gửi tới \textbf{Amazon SageMaker} để
phân loại mức độ nguy cấp và tổng hợp báo cáo y khoa. 2. \textbf{Hệ
thống đặt lịch chống trùng slot:} Sử dụng cơ chế khóa dòng
(\textbf{Pessimistic Locking}) trên \textbf{AWS RDS PostgreSQL} để giải
quyết triệt để bài toán Race Condition khi nhiều người cùng đặt một
khung giờ. 3. \textbf{Thanh toán linh hoạt 2 đợt:} Hỗ trợ thanh toán
online đợt 1 (tiền khám gốc qua PayOS/VNPay/Stripe) và thanh toán đợt 2
(tiền thuốc/xét nghiệm phát sinh tại quầy). 4. \textbf{Bảo mật và Phân
quyền (RBAC):} Quản lý định danh qua \textbf{AWS Cognito} và bảo mật dữ
liệu y tế nhạy cảm bằng mã hóa \textbf{UUID}.

\emph{Lợi ích và hoàn vốn đầu tư (ROI)}\\
- Giảm đến 60\% thời gian chờ đợi tại quầy nhờ quy trình check-in tự
động bằng mã QR. - Tăng 30\% hiệu suất làm việc của Bác sĩ nhờ có sẵn
báo cáo tóm tắt triệu chứng do AI tổng hợp. - Triệt tiêu 100\% rủi ro
trùng lịch khám, nâng cao trải nghiệm và sự hài lòng của bệnh nhân. -
Tối ưu hóa chi phí hạ tầng nhờ khả năng linh hoạt của AWS Cloud theo lưu
lượng truy cập thực tế.

\vspace{0.5em}\noindent\textbf{3. Kiến trúc giải pháp}

Hệ thống sử dụng kiến trúc Web Multi-tier trên nền tảng AWS Cloud, chia
làm 3 lớp chính: Frontend (Next.js), Backend Service (NestJS trên EC2),
và AI Services (Amazon SageMaker).

\emph{Dịch vụ AWS sử dụng}\\
- \emph{Amazon Cognito}: Quản lý định danh, đăng nhập/đăng ký và phân
quyền người dùng (Patient, Doctor, Receptionist, Admin). - \emph{AWS
EC2}: Triển khai ứng dụng Web Frontend (Next.js) và Backend REST API /
WebSocket (NestJS) qua Docker Containers. - \emph{Amazon SageMaker}:
Endpoint lưu trữ và chạy mô hình AI phân tích triệu chứng y tế. -
\emph{Amazon RDS (PostgreSQL)}: Cơ sở dữ liệu quan hệ lưu trữ thông tin
bệnh nhân, bác sĩ, lịch làm việc và hóa đơn. - \emph{Amazon S3}: Lưu trữ
file tĩnh, ảnh đại diện, đơn thuốc và kết quả cận lâm sàng. - \emph{AWS
CloudWatch}: Thu thập log hệ thống, giám sát hiệu năng CPU/Memory và gửi
cảnh báo sự cố. - \emph{AWS SES / SNS}: Tự động gửi Email/SMS xác nhận
lịch hẹn kèm mã QR check-in.

\emph{Thiết kế thành phần}\\
- \emph{Giao diện Bệnh nhân (Next.js)}: Cho phép chat với Bot, xem danh
sách bác sĩ, chọn khung giờ và thanh toán online. - \emph{Giao diện Nhân
viên Y tế (Next.js)}:\\
- \textbf{Doctor Portal:} Xem lịch hẹn trong ngày, đọc báo cáo tóm tắt
AI, nhập kết quả chẩn đoán và đơn thuốc. - \textbf{Reception Portal:}
Check-in bệnh nhân qua mã QR, xuất hóa đơn và thu tiền đợt 2 tại quầy. -
\emph{Backend Microservices (NestJS)}: Xử lý Business Logic, WebSocket
Gateway, tích hợp Cổng thanh toán (PayOS/VNPay/Stripe) và kết nối CSDL
RDS PostgreSQL.

\vspace{0.5em}\noindent\textbf{4. Triển khai kỹ thuật}

\emph{Các giai đoạn triển khai}\\
Dự án được triển khai trong vòng \textbf{8 tuần} (2 tháng) qua các giai
đoạn:\\
1. \emph{Tuần 1 - 2}: Phân tích yêu cầu, thiết kế kiến trúc Cloud AWS,
thiết kế CSDL và khởi tạo Base Source Code (Next.js + NestJS). 2.
\emph{Tuần 3 - 4}: Triển khai hạ tầng AWS (EC2, S3, RDS PostgreSQL),
tích hợp ORM Prisma/TypeORM và xử lý cơ chế chống trùng lịch. 3.
\emph{Tuần 5 - 6}: Tích hợp AWS CloudWatch, tích hợp mô hình AI trên
Amazon SageMaker, hoàn thiện cổng thanh toán (PayOS/VNPay/Stripe) và
Dashboard phân quyền (RBAC). 4. \emph{Tuần 7 - 8}: Kiểm thử tải (Stress
Test), kiểm thử E2E, tối ưu Indexing CSDL, đóng gói tài liệu Swagger API
và tổng kết Demo.

\emph{Yêu cầu kỹ thuật}\\
- \emph{Frontend}: Next.js 14, TailwindCSS, Socket.io-client, React
Query. - \emph{Backend}: NestJS Framework, TypeORM/Prisma, TypeScript,
Socket.io. - \emph{Database}: PostgreSQL 16.x trên AWS RDS (Private
Subnet). - \emph{AI/ML}: Python, Amazon SageMaker Endpoints. -
\emph{DevOps}: Docker, AWS CLI, GitHub Actions (CI/CD).

\vspace{0.5em}\noindent\textbf{5. Lộ trình \& Mốc triển khai}

\begin{itemize}
\tightlist
\item
  \emph{Giai đoạn 1 (Tuần 1 - Tuần 2)}: Hoàn thiện thiết kế System
  Architecture \& CSDL Schema.
\item
  \emph{Giai đoạn 2 (Tuần 3 - Tuần 4)}: Triển khai thành công hạ tầng
  AWS (RDS PostgreSQL, S3, EC2) \& API Core.
\item
  \emph{Giai đoạn 3 (Tuần 5 - Tuần 6)}: Tích hợp thành công SageMaker
  AI, Payment Gateway (PayOS/VNPay/Stripe) \& CloudWatch.
\item
  \emph{Giai đoạn 4 (Tuần 7 - Tuần 8)}: Hoàn thành Stress Test, tối ưu
  UX/UI và nghiệm thu dự án.
\end{itemize}

\vspace{0.5em}\noindent\textbf{6. Ước tính ngân sách}

Chi phí hạ tầng được ước tính trên môi trường AWS Cloud cho giai đoạn
thử nghiệm (MVP / Development Phase):

\emph{Chi phí hạ tầng AWS hàng tháng}\\
- \emph{AWS RDS (db.t3.micro - PostgreSQL)}: \textasciitilde15.00
USD/tháng (Single-AZ, 20 GB Storage). - \emph{AWS EC2 (t3.small -
Backend \& Frontend)}: \textasciitilde14.00 USD/tháng (Chạy 24/7). -
\emph{Amazon SageMaker (ml.t3.medium Endpoint)}: \textasciitilde36.00
USD/tháng (Phục vụ Serverless/On-demand AI Inference). - \emph{Amazon S3
Standard}: \textasciitilde0.50 USD/tháng (5 GB Data Storage \&
Transfer). - \emph{Amazon Cognito}: 0.00 USD/tháng (Miễn phí 50,000 MAU
đầu tiên). - \emph{AWS CloudWatch}: \textasciitilde2.00 USD/tháng
(Metrics, Logs \& Alarms). - \emph{Cổng thanh toán (PayOS / VNPay /
Stripe Sandbox)}: 0.00 USD (Môi trường Test miễn phí).

\emph{Tổng chi phí hạ tầng Cloud}: \textasciitilde67.50 USD/tháng.

\vspace{0.5em}\noindent\textbf{7. Đánh giá rủi ro}

\emph{Ma trận rủi ro}\\
- Trùng lịch khám do thao tác đồng thời (Race Condition): Ảnh hưởng cao,
xác suất trung bình. - Gián đoạn kết nối mô hình AI (SageMaker Timeout):
Ảnh hưởng trung bình, xác suất thấp. - Lộ thông tin y tế bệnh nhân (Data
Privacy): Ảnh hưởng rất cao, xác suất thấp.

\emph{Chiến lược giảm thiểu}\\
- \emph{Race Condition}: Sử dụng Pessimistic Locking trực tiếp từ tầng
Database PostgreSQL.\\
- \emph{AI Timeout}: Xây dựng luồng Fallback --- Nếu AI bị lỗi, hệ thống
tự động chuyển sang luồng chọn bác sĩ thủ công truyền thống mà không làm
gián đoạn trải nghiệm người dùng.\\
- \emph{Data Privacy}: Sử dụng chuỗi UUID thay cho ID tự tăng trên URL,
ẩn danh thông tin bệnh nhân khi lưu log lên CloudWatch.

\vspace{0.5em}\noindent\textbf{8. Kết quả kỳ vọng}

\begin{itemize}
\tightlist
\item
  \textbf{Cải tiến kỹ thuật:} Tự động hóa 80\% quy trình tiếp nhận và
  phân luồng bệnh nhân; đạt chỉ số sẵn sàng của hệ thống \textgreater{}
  99.9\% trên nền tảng AWS Cloud.\\
\item
  \textbf{Giá trị dài hạn:} Cung cấp nguồn dữ liệu chuẩn hóa cho công
  tác phân tích và dự báo dịch bệnh; hệ thống dễ dàng mở rộng cho chuỗi
  nhiều chi nhánh phòng khám trong tương lai.
\end{itemize}
